\par\noindent\textbf{(1)} 平面运动单位质量能量E=(Ṙ²+R²φ̇²)/2+Φ(R)，角动量L=R²φ̇。消去φ̇得到径向有效势V\_eff(R)=Φ(R)+L²/(2R²)。圆轨道须V\_eff'=Φ'−L²/R³=0，故L²=R0³Φ'(R0)=R0⁴Ω0²。
\par\noindent\textbf{(2)} 固定L在线性近似R=R0+ξ，ξ̈=−V\_eff''(R0)ξ，于是κ²=Φ''(R0)+3Φ'(R0)/R0。由Ω²(R)=Φ'(R)/R可化为κ²=R dΩ²/dR+4Ω²，这也是轴对称星系局部环振动频率。
\par\noindent\textbf{(3)} Kepler势Φ=−GM/R给Ω²=GM/R³，因R dΩ²/dR=−3Ω²，故κ=Ω；谐振势Φ=ω²R²/2给Ω=ω、κ=2ω；平坦旋转曲线v\_c常量时Ω=v\_c/R，κ=√2Ω。三例显示相同圆速不能决定径向频率，关键是旋转曲线梯度。
\par\noindent\textbf{(4)} 若κ²<0，扰动ξ∝e\textasciicircum{}\{|κ|t\}指数增长，圆轨道不是局部稳定极小值；若κ²>0，径向扰动谐振且能量二阶增量δE=[ξ̇²+κ²ξ²]/2守恒。κ→0是边缘稳定，非线性项及外扰动可能主导。
